\begin{examplebox}
\paragraph{例1：先看导数符号，再分别写区间}
求 $f(x)=x^3-3x+1$ 的单调区间。
\[
f'(x)=3(x-1)(x+1).
\]
以 $-1,1$ 分段，导数符号依次为 $+,-,+$，所以严格递增区间可写为 $(-\infty,-1]$、$[1,+\infty)$，严格递减区间为 $[-1,1]$。端点可纳入，因为函数连续；各段必须分别陈述。
\begin{center}
\begin{tikzpicture}
\begin{axis}[width=11.6cm,height=6.2cm,axis lines=middle,
 xmin=-2.2,xmax=2.2,ymin=-3.3,ymax=5.3,
 xlabel={$x$},ylabel={$y$},xtick=\empty,ytick={-1,3},
 title={$y=x^3-3x+1$：两段递增，中间递减},title style={font=\small}]
\addplot[blue,very thick,domain=-2.1:-1,samples=60]{x^3-3*x+1};
\addplot[orange!85!black,very thick,dashed,domain=-1:1,samples=60]{x^3-3*x+1};
\addplot[blue,very thick,domain=1:2.1,samples=60]{x^3-3*x+1};
\addplot[only marks,black] coordinates {(-1,3) (1,-1)};
\node[font=\small,above] at (axis cs:-1,3) {$(-1,3)$};
\node[font=\small,below] at (axis cs:1,-1) {$(1,-1)$};
\end{axis}
\end{tikzpicture}
\end{center}

回验：$f(-1)=3$，$f(1)=-1$，与中间一段下降一致。

\paragraph{例2：导数为零不一定破坏严格递增}
证明 $g(x)=x^3$ 在 $\mathbb R$ 上严格递增。
最低门槛是直接作差：任取 $u<v$，
\[
v^3-u^3=(v-u)(u^2+uv+v^2)>0,
\]
其中 $u^2+uv+v^2=(u+v/2)^2+3v^2/4>0$（两数不能同时为零）。这也回验了 $g'(0)=0$ 与严格递增并不矛盾。
\end{examplebox}
